\hypertarget{uxc5f0uxad6c-uxb85cuxadf8-uxace0uxcc28uxc6d0-uxad6c-uxcc44uxc6b0uxae30sphere-packing-uxd558uxd55c-uxac1cuxc120-uxc2dcuxb3c4}{%
\subsection*{연구 로그: 고차원 구 채우기(sphere packing) 하한 개선 시도}\label{uxc5f0uxad6c-uxb85cuxadf8-uxace0uxcc28uxc6d0-uxad6c-uxcc44uxc6b0uxae30sphere-packing-uxd558uxd55c-uxac1cuxc120-uxc2dcuxb3c4}}

\begin{itemize}
\tightlist
\item
  \textbf{문제}: \(\delta_n\) = \(\mathbb{R}^n\) 격자 구 채우기 최대 밀도. 알려진 하한 \(\delta_n \ge c\,n^2 2^{-n}\) (Klartag 2025, arXiv:2504.05042) 및 그 후속(\(\log\log n\) 인자 추가)을 개선하거나, 개선의 장애물을 정확히 규명한다.
\item
  \textbf{시작 시각}: 2026-11-25 06:24:15 UTC (epoch 1795587855). \textbf{종료 조건}: 5시간 경과(11:24:15 UTC) 또는 30회차 완료.
\item
  \textbf{표기}: \(V\) := (covolume 1 격자에 대해 원점 외 격자점이 없는 열린 공의 부피) \(= 2^n\delta\). Minkowski--Hlawka: \(V\ge 2\zeta(n)\); Rogers/Ball: \(V\gtrsim n\); Venkatesh: \(V\gtrsim n\log\log n\) (희소 차원열); Klartag: \(V\ge cn^2\). 상한: Minkowski \(V\le 2^n\), KL \(V\le 2^{0.401n}\), (업로드된 OpenAI 논문) \(V\le 2^{(1-0.6044+o(1))n}\).
\item
  \textbf{분류 규칙}: {[}증명됨{]} / {[}조건부{]} / {[}추측{]} / {[}수치관찰{]} / {[}휴리스틱{]} / {[}실패{]} 로 태그한다.
\end{itemize}

\begin{center}\rule{0.5\linewidth}{0.5pt}\end{center}

\hypertarget{uxd68cuxcc28-uxc900uxbe44-06240650-utc-uxb450-uxb17cuxbb38-uxc815uxb3c5-uxbc0f-uxcd08uxae30-uxbd84uxc11d}{%
\subsubsection*{0회차 (준비, 06:24--06:50 UTC): 두 논문 정독 및 초기 분석}\label{uxd68cuxcc28-uxc900uxbe44-06240650-utc-uxb450-uxb17cuxbb38-uxc815uxb3c5-uxbc0f-uxcd08uxae30-uxbd84uxc11d}}

\hypertarget{uxc77duxc740-uxb0b4uxc6a9-uxc694uxc57d}{%
\paragraph{읽은 내용 요약}\label{uxc77duxc740-uxb0b4uxc6a9-uxc694uxc57d}}

\begin{enumerate}
\def\labelenumi{\arabic{enumi}.}
\tightlist
\item
  \textbf{Klartag (하한)}: 랜덤 격자 \(L\) (covol \(=\mathrm{Vol}(B^n)\))과 반지름 \(\approx 1-1/n\) 공에서 출발, 대칭행렬 공간에서 Dyson 브라운 운동 \(dA_t=\pi_t(dW_t)\) 로 타원체 \(\mathcal E_t=\{A_tx\cdot x<1\}\) 를 진화. 격자점에 닿으면 그 점을 경계에 고정(선형제약 \(A x\cdot x=1\)), 제약공간 \(F_t\) 의 차원 \(N_t\) 가 0이 되면 정지. Itô 공식: \(\mathbb E\log\det A_T\le n\log a_0-\tfrac12\int_0^T\mathbb E[\|A_t\|^{-2}N_t]dt\). 접촉점 기대개수는 Siegel 평균 + 반사원리로 \(\lesssim e^{n^2t/8}/n^2\) 적분 유계 ⇒ \(T=16\log n/n^2\) 에서 \(\det A_T\le C/n^4\), 즉 부피 \(\gtrsim n^2\).

  \begin{itemize}
  \tightlist
  \item
    핵심 수: 대칭행렬 공간 차원 \(n(n+1)/2\) = 타원체 자유도 = 흡수할 수 있는 접촉점 쌍의 수. "쓸어 담은 부피 ≈ 만난 격자점 수(Poisson)" ⇒ \(V\sim n^2\).
  \end{itemize}
\item
  \textbf{OpenAI "Ten Advances" Ch.1--2 (상한)}: Cohn--Elkies LP 의 지수적 증가율이 정확히 \(\sqrt{e/(2\pi)}\), 즉 \(\Delta_d\le 2^{-(0.6044+o(1))d}\) (KL 지수 0.5990 개선). LP 방법 자체의 한계도 같은 지수.
\end{enumerate}

\hypertarget{uxcd08uxae30-uxad00uxcc30-uxac80uxc99d-uxc804-uxbcf8-uxc138uxc158uxc5d0uxc11c-uxac80uxd1a0uxd560-uxac83}{%
\paragraph{초기 관찰 (검증 전, 본 세션에서 검토할 것)}\label{uxcd08uxae30-uxad00uxcc30-uxac80uxc99d-uxc804-uxbcf8-uxc138uxc158uxc5d0uxc11c-uxac80uxd1a0uxd560-uxac83}}

\begin{itemize}
\tightlist
\item
  \textbf{{[}휴리스틱 H1{]} Klartag 과정 = Ryshkov 다면체 위 "끈끈한(sticky)" 브라운 운동.} \(L\)-free 행렬 집합 \(\mathcal R_L=\{A\succ0: Ax\cdot x\ge1\ \forall x\in L\setminus0\}\) 는 Ryshkov 다면체, 꼭짓점 = perfect form. Klartag 과정은 면(face)을 따라 차원이 줄어들며 꼭짓점(perfect form)에서 멈춘다. 따라서 "명시적 구성" 관점의 고전적 대응물은 Voronoi 이론(perfect/eutactic form)이다.
\item
  \textbf{{[}휴리스틱 H2{]} 대칭 격자족의 이득 공식.} 유한군 \(G\) 가 \(\mathbb Q\)-기약 표현 \(W\) (\(\mathbb QG\) 의 단순성분 \(A=M_r(\Delta)\)) 의 \(k\) 겹 \(W^k\) 에 작용, \(n=k\dim_{\mathbb Q}W\). \(G\)-불변 이차형식 공간 차원 \(D\approx k^2\dim\Delta/2\), 일반 궤도 크기 \(|G|\) ⇒ Klartag×G 휴리스틱 \(V\sim D|G|\sim n^2\cdot|G|/\dim_{\mathbb Q}A\). 또한 \(|G|/\dim_{\mathbb Q}A=\sum_{g\in G}\rho(g)^2\), \(\rho=\chi_W/\chi_W(1)\) (검증 필요). 순환군 \(\mu_m\): \(m/\varphi(m)\sim e^\gamma\log\log m\) ⇒ "\(n^2\log\log n\)" 과 일치할 것으로 추정.
\item
  \textbf{{[}휴리스틱 H3{]}} 고정점 없는(fpf) 작용만 쓰면 이득은 \(O(\log\log)\) 로 보임(Frobenius complement 구조). 더 큰 이득은 non-fpf 군(Heisenberg, \(SL_2(p)\) Weil, Clifford 군 등)이 필요하나, 고정 부분격자 \(\Lambda\cap U^k\) 가 짧은 벡터를 만들지 않으려면 "초안정성" \(\mathrm{covol}(W_\mathbb Z\cap U)^{2/\dim U}\dim U\gtrsim \mathrm{covol}(W_\mathbb Z)^{2/\dim W}\dim W\) (모든 \(\Delta\)-부분공간 \(U\)) 가 필요해 보임. \(r\ge2\) 이면 \(\Delta\) 의 root discriminant \(\gtrsim 2\pi e/c_0\) 가 필요조건으로 보임.
\item
  \textbf{{[}휴리스틱 H4{]} 구멍 확률 관점.} \(P_V:=\Pr_{\rm Haar}(L\cap B_V=\{0\})\) 의 지지 끝점이 \(2^n\delta_n\). Poisson 휴리스틱 \(P_V\approx e^{-V/2}\) 과 "최적 격자의 근방 측도 \(\ge e^{-Cn^2\log n}\)" 를 비교하면 \(V\sim n^2\log n\) 이 자연스러운 장벽으로 보임.
\end{itemize}

\hypertarget{uxb2e4uxc74c-uxd68cuxcc28-uxacc4uxd68d}{%
\paragraph{다음 회차 계획}\label{uxb2e4uxc74c-uxd68cuxcc28-uxacc4uxd68d}}

\begin{itemize}
\tightlist
\item
  \textbf{1회차}: 문헌 조사 --- 최신 \(\log\log n\) 결과의 정확한 출처/정리/방법, 대칭 격자(Venkatesh, Vance, Gargava, Gargava--Viazovska), random perfect lattice 수치연구(Andreanov--Scardicchio, Kallus), 구멍확률/Poisson 통계.
\item
  \textbf{2회차}: 소차원 수치실험 --- Klartag 과정 시뮬레이션(\(n\le 10\) 내외), 알려진 최밀 격자의 \(V_n\) 과 \(n^2\) 비교표, 접촉점(키싱 수) 패턴.
\item
  \textbf{3회차}: H2--H3 의 엄밀화 --- 대칭 Klartag 정리의 정확한 서술·증명 골격, non-fpf 장애물/가능성 판정.
\end{itemize}

\begin{center}\rule{0.5\linewidth}{0.5pt}\end{center}

\hypertarget{uxd68cuxcc28-06480711-utc-uxacbduxacfc-47uxbd84-uxbb38uxd5cc-uxc870uxc0ac-uxcd5cuxc2e0-uxacb0uxacfcuxc640-uxc54cuxb824uxc9c4-uxc7a5uxbcbd}{%
\subsubsection*{1회차 (06:48--07:11 UTC, 경과 47분): 문헌 조사 --- 최신 결과와 알려진 장벽}\label{uxd68cuxcc28-06480711-utc-uxacbduxacfc-47uxbd84-uxbb38uxd5cc-uxc870uxc0ac-uxcd5cuxc2e0-uxacb0uxacfcuxc640-uxc54cuxb824uxc9c4-uxc7a5uxbcbd}}

\textbf{선택한 과제}: "log log n 이 붙은 최신 결과"의 정확한 내용·방법과, 대칭 격자족/perfect form/Poisson 통계에 대해 알려진 사실 확정. (에이전트 4개 병렬; 노트: \texttt{rounds/R1\_lit\_*/notes.md}, 논문: \texttt{papers/})

\textbf{확정된 사실 {[}문헌{]}}
- \textbf{현재 최선의 하한}: Gargava--Gröbner--Simkin (arXiv:2610.12619, 2026-10): 모든 큰 \(d\) 에 대해 \(V_d\ge c\,d^2\log\log d\). 그 직전 Abuya--Gargava--Zhao (arXiv:2606.05105, 2026-06): 희소 차원열 \(N=\varphi(m)r\), \(r=\lceil(\log\varphi(m))^4\rceil\) 에서 같은 한계. 둘 다 "Klartag 과정을 원분체 대칭(\(\mu_m\))-불변 이차형식 공간 \(\bigoplus \mathrm{Herm}_r(\mathbb C)\) 안에서 실행": \textbf{궤도(크기 \(m\)) 하나가 선형제약 1개만 소모} ⇒ 부피 ≈ (궤도 크기)×(불변형식 차원) \(= m\cdot Nr/2 = N^2\,m/\varphi(m)\). 랭크 \(r\) 을 키워야 하는 이유: 고정 \(r\) 이면 접촉 수가 \(V^{1+1/r}\) 로 늘어 \((mNr)^{r/(r+1)}\) 밖에 못 얻음 (AGZ Remark 4.3).
- \textbf{상수}: Rezvyakova (arXiv:2608.22590): Klartag 정리의 상수 \(c=\tfrac1{2e}-o(1)\), 즉 \(V_n\ge(1-o(1))N/e\), \(N=n(n+1)/2\).
- \textbf{Poisson 근사의 한계}: Klartag (arXiv:2605.10479): 부피 \(\le cn\) 인 집합에서만 Haar-랜덤 격자 ≈ Poisson (TV 거리 \(\le Ce^{-c'n}\)); 부피 \(\ge n\) 에서는 일반적으로 거짓. Schmidt/Kim: 구멍확률 \(P_V=e^{-V/2}+O(e^{-cn})\) 은 \(V<Cn\) 에서만. \(Cn\le V\le cn^2\) 구간에서 \(P_V\) 에 대해 알려진 것은 \(e^{-(1+o(1))n^2\log n}\le P_V\le e^{-cn}\) 뿐.
- \textbf{장벽에 관한 유일한 명시적 진술}: Venkatesh 의 추측 \(V_N\le N^2(\log N)^{O(1)}\) ("격자의 자유도가 \(N^2\)").
- \textbf{대칭 방법의 한계 {[}문헌{]}}: Gargava(2021): 나눗셈대수의 유한부분군 \(G_0\) 로는 \(|G_0|/\dim_\mathbb Q D\le 3m/\varphi(m)\) 이라 \(O(d\log\log d)\) 를 넘지 못함.
- \textbf{물리 문헌의 소차원 데이터}: Kallus--Marcotte--Torquato(2013): 순차 LP 로 얻은 isostatic extreme 격자의 전형적 밀도 \(2^d\langle\phi\rangle\approx0.0567\,d^{3.01}\) (\(d=15..24\)); Andreanov--Scardicchio(2012): 랜덤 perfect 격자 \(V_{typ}/d^2\in[0.65,0.94]\) (\(d=8..19\); "\(2^{-0.84d}\)" 는 구분 불가능한 두 fit 중 하나일 뿐). 9차원 perfect form 2,237,251,040 개 완전 분류(Dutour Sikirić--van Woerden 2025): 60.5\%가 최소 접촉수 45쌍.
- 찾지 못한 것: \(n^2\log n\) 이상의 결과, Klartag 과정의 수치 시뮬레이션, 명시적(비확률적) 구성, \(N^2\mathrm{polylog}\) 상계의 증명.

\textbf{다음 할 일}: (a) 소차원 시뮬레이션으로 Klartag 과정의 실제 상수 확인(2회차). (b) 대칭 이득 공식과 비자유 작용 분석(3회차).

\begin{center}\rule{0.5\linewidth}{0.5pt}\end{center}

\hypertarget{uxd68cuxcc28-06480719-utc-uxacbduxacfc-55uxbd84-uxc18cuxcc28uxc6d0-uxc218uxce58uxc2e4uxd5d8-n2..10}{%
\subsubsection*{\texorpdfstring{2회차 (06:48--07:19 UTC, 경과 55분): 소차원 수치실험 (\(n=2..10\))}{2회차 (06:48--07:19 UTC, 경과 55분): 소차원 수치실험 (n=2..10)}}\label{uxd68cuxcc28-06480719-utc-uxacbduxacfc-55uxbd84-uxc18cuxcc28uxc6d0-uxc218uxce58uxc2e4uxd5d8-n2..10}}

\textbf{선택한 과제}: 사용자가 제안한 "작은 \(n\) 에서 실제로 해 보기". (i) Klartag 과정을 동결될 때까지 실행, (ii) Ryshkov 다면체 위 sticky walk + 탐욕 LP 하강으로 extreme form 표본 추출. 각각 독립 검증 에이전트가 재현·검산. (코드/데이터: \texttt{rounds/R2\_klartag\_sim/}, \texttt{rounds/R2\_voronoi\_descent/})

\textbf{{[}수치관찰{]} (검증자가 독립 열거 코드로 전부 재확인; 버그 없음)}
- 검산: \(n=2\) 항상 육각격자, \(n=3\) 항상 FCC, \(n=4\) 는 \(D_4/A_4\) 만, \(n=5\) 는 알려진 perfect form 3개만, \(n=8\) 에서 \(E_8\)(120쌍) 재현. 저장된 최종 form 3403+1878 개 모두 격자점 없음·perfect; 어떤 \(V\) 도 알려진 최적값 초과 없음.
- \textbf{Klartag 과정의 최종 부피는 \(V_{final}\approx N=n(n+1)/2\)}: \(\mathbb E[V_{final}]/N=1.21,0.99,0.93,0.97,1.08,1.26,1.24,1.01,1.03\) (\(n=2..10\)). 경로 전체에서 \(V(t)/k(t)\approx1\) (\(k\)=독립 접촉쌍 수; \(k\ge36\) 에서 0.96--1.02). 즉 "쓸어 담은 부피 = 접촉쌍 수" 가 상수 1로 성립 ⇒ 실제 상수는 \(c\approx1/2\) (증명된 값은 \(1/(2e)\)).
- 과정이 최밀 격자에서 끝날 확률: \(n=2..8\): 1, 1, 0.47, 0.85, 0.67, 0.79, 0.33; \(n=9,10\): 0/100, 0/50. 즉 \(n\le8\) 에서는 root 격자가 지배하는 "작은 차원 효과"이고 \(n\ge9\) 부터 \(n^2\) 법칙으로 떨어진다(평균/최선 = 0.69, 0.61, 0.60).
- 탐욕 하강 후 extreme form: 중앙값 \(V\approx(0.63..0.70)n^2\) (\(n=9,10\)), 접촉쌍 평균 \(\approx2N\), 최선 격자는 \(3.0\)--\(3.3N\). \(V\) 와 접촉쌍 수의 Spearman 상관 0.88--0.95.
- \textbf{구조}: 발견된 모든 extreme form 은 작은 분모의 유리 형식(\(qA\) 정수, \(q\in\{2,4,..,14\}\); root 격자만 \(q=2\), 대부분 \(q=4\): 최소 노름 4인 정수격자, 최소벡터 내적 \(\in\{0,\pm1,\pm2\}\)).
- \(n\le10\) 데이터로는 다항/지수 성장 구분 불가(검증자: fit 지수는 의미 없음).

\textbf{해석 {[}휴리스틱{]}}: 작은 \(n\) 의 "규칙성"(root 격자, laminated 격자)은 점근 영역(\(V\sim n^2\))과 다른 체제다. 알려진 최밀 격자는 \(n\le 72\) 에서 \(V/n^2\) 가 지수적으로 큼(\(n=24\): 56, \(n=48\): 2832). 따라서 소차원 패턴의 외삽으로 점근 하한을 얻기는 어렵고, 점근 영역의 구조는 "접촉쌍 수 ↔ 부피" 회계로 봐야 한다.

\textbf{실패/주의}: Goldstein--Mayer 격자(\(p=100003\))는 \(n\ge8\) 에서 Haar 의 거친 근사(검증자 지적). \(n\ge11\) 미수행.

\begin{center}\rule{0.5\linewidth}{0.5pt}\end{center}

\hypertarget{uxd68cuxcc28-06480741-utc-uxacbduxacfc-77uxbd84-uxb300uxce6d-uxaca9uxc790uxc871-uxc774uxb860-uxc99duxba85-uxc5d0uxc774uxc804uxd2b8-3-uxbc18uxbc15-uxc5d0uxc774uxc804uxd2b8}{%
\subsubsection*{3회차 (06:48--07:41 UTC, 경과 77분): 대칭 격자족 이론 (증명 에이전트 3 + 반박 에이전트)}\label{uxd68cuxcc28-06480741-utc-uxacbduxacfc-77uxbd84-uxb300uxce6d-uxaca9uxc790uxc871-uxc774uxb860-uxc99duxba85-uxc5d0uxc774uxc804uxd2b8-3-uxbc18uxbc15-uxc5d0uxc774uxc804uxd2b8}}

\textbf{선택한 과제}: 0회차의 H2--H3 를 엄밀화. (노트: \texttt{rounds/R3\_G\_klartag/notes.md}, \texttt{rounds/R3\_non\_free/notes.md}; 반박 보고는 뒤 회차에 반영)

\textbf{{[}증명됨{]} (노트에 증명 수록, 수치 검산)}
- \textbf{이득 항등식}: \(W\) 가 \(\mathbb Q\)-기약, \(A=\mathbb Q G\text{-span}=M_r(\Delta)\) 이면 \(\sum_{g\in G}\rho(g)^2=|G|/\dim_\mathbb QA\), \(\rho=\chi_W/\chi_W(1)\). 불변형식 차원 \(D=\frac1{|G|}\sum_g\frac{\chi(g)^2+\chi(g^2)}2\approx k^2\dim_\mathbb Q\Delta/2\). \textbf{Klartag 대비 이득은 \(g=|G|/(2\dim_\mathbb QA)\)} (0회차 H2 의 인자 2 수정). 고정점 없는 작용이면 \(\rho(g)=\mu(o(g))/\varphi(o(g))\) 이므로 \(|G|/\dim A=\sum_{C\le G\text{ 순환, 제곱인수 없는 위수}}1/\varphi(|C|)\).
- \textbf{고정점 없는 군의 상한}: Frobenius complement 의 구조정리(Zassenhaus)로 \(g\le 24\,c/\varphi(c)=O(\log\log|G|)\); 순환군이 \(e^\gamma\log\log m\) 달성. ⇒ \textbf{자유 작용에 기반한 Venkatesh/AGZ 형 논증은 \(\log\log\) 를 넘을 수 없다.} (상수는 \(SL_2(5)\times C_N\) 이 4배.)
- \textbf{AGZ 구조의 독립 재유도}: 사영된 이차변동 \(q(x)=\frac12\sum_\sigma u_\sigma^2\), \(\mathbb E q=(k+1)/(n+2)\); 접촉 수 \(\sim V^{1+1/k}\); 희소 벡터 때문에 \(k>0.407\log(\cdot)\) 필요. (AGZ 논문과 일치함을 1회차 문헌으로 확인.)
- \textbf{비자유 작용의 평균값 공식}: \(k>r\) 일 때 \(\mathbb E\#(\Lambda\cap B)=V+\sum_{U\subsetneq W}T_U\), \(T_U=\mathrm{Vol}_{fn}(R)/c_U^k\approx V^f\mathrm{poly}(n)\,R_U^{k/2}\), \(R_U=(1/f)^{\dim U}c_W^{2\dim U/\dim W}/c_U^2\). \textbf{텐서 안정성}(\(R_U<1\ \forall U\))이 필요충분(지수 스케일).
- \textbf{No-go 정리들}: (1) \(\Delta=\mathbb Q\), \(r\ge2\) 이면 항상 실패(\(\gamma_r/r\le2/\pi<1\)) --- Weyl 군, \(S_n\), 실 Clifford 군/Barnes--Wall 전부 배제. (2) 직선 조건 ⇒ \(\mathrm{rd}_\Delta\ge rd/\gamma_{rd}\to2\pi e/1.744\). (3) 소수 위수 \(q\) 원소가 \(W|_{\langle g\rangle}=a\cdot\mathbf1+b\cdot\mathbb Q(\zeta_q)\) 이면 \((a+b(q-1))\,h\!\big(\tfrac a{a+b(q-1)}\big)\le2\min(a,b)\log q\); \(q=2\) 면 \(a=b\) 이고 정확히 경계(\(R_U=1\)). (4) \(f\tau_U^2<0.5735\) 면 어떤 \(L\) 로도 지수적 손실.
- 반례/정정: "자유 작용 ⇒ \(A\) 가 나눗셈대수"는 거짓(\(SL_2(5)\times C_7\): \(A=M_2(K)\)).

\textbf{{[}미해결 후보{]}}: \(SL_2(p)\), \(p\equiv1\ (4)\), 충실한 Weil 표현(차수 \((p-1)/2\), \(\Delta\)=\(\mathbb Q(\sqrt p)\) 위 정부호 사원수대수, \(r=(p-1)/4\), \(\dim_\mathbb QW=2(p-1)\); \(p=5\): \(E_8\), \(p=13\): Leech). 유도한 모든 필요조건을 통과. 이득 \(g\sim p\) ⇒ (성립한다면) \(V\sim n^{2.5}\) --- Venkatesh 추측과 충돌하므로 큰 \(p\) 에서 장애물이 있을 것으로 예상. → 5회차에서 검토.

\textbf{{[}실패{]} 로 판정된 후보들}: \(PSL_2(p)\) (\(p\equiv3\) mod 4) Craig 격자: 직선에서 \(f\tau^2\to1/2\); Heisenberg/표준격자 \(O_K^p\): \(1/p\); Barnes--Wall; \(A_2\otimes L\) 등.

\hypertarget{uxd68cuxcc28-uxc911-uxbcf8uxc778uxc624uxcf00uxc2a4uxd2b8uxb808uxc774uxd130-uxc720uxb3c4-45uxd68cuxcc28uxc758-uxcd9cuxbc1cuxc810}{%
\paragraph{3회차 중 본인(오케스트레이터) 유도 --- 4·5회차의 출발점}\label{uxd68cuxcc28-uxc911-uxbcf8uxc778uxc624uxcf00uxc2a4uxd2b8uxb808uxc774uxd130-uxc720uxb3c4-45uxd68cuxcc28uxc758-uxcd9cuxbc1cuxc810}}

\begin{itemize}
\tightlist
\item
  \textbf{{[}휴리스틱→4회차{]} 퍼셉트론 모형.} Poisson 가정에서 "부피 \(v\) 인 공을 \(\|B_0\|_{HS}=b\) 만큼 변형한 타원체가 격자점을 피할 조건"은 \(\langle\theta,g_i\rangle\ge-\sigma_i\) (\(\theta=\sqrt2B_0\in\mathbb R^{N-1}\), \(g_i\) 가우스, \(\sigma_i\) 는 강도 \(\frac v4e^{\sigma/2}d\sigma\) 의 Poisson 과정) 꼴의 \textbf{랜덤 선형계획/음의 마진 퍼셉트론} 문제. Gordon 비교로 볼록 부분의 임계값 \(v^*=(2/e)N\), 그때 \(\|\theta\|^2=8\), 활성 제약 \(N\) 개, 타원체 부피 \(V=e\,v^*=2N\approx n^2\). 즉 (모형에서) 대각합-LP 최적 상수는 1, Klartag 과정은 \(1/2\), 증명된 값은 \(1/(2e)\). 비볼록 영역(큰 변형 \(\rho\))에서는 1차 모멘트 상계 \(V\lesssim N\log\rho\); 존재성은 2차 모멘트가 필요 --- 음의 퍼셉트론(Montanari--Zhong--Zhou)의 "존재 임계 vs 알고리즘 임계 사이 \(\log|\kappa|\) 간극"과 같은 구조. ⇒ "알고리즘적(온라인) 방법은 \(\Theta(n^2)\), 비구성적 존재는 \(n^2\log n\) 까지"라는 작업가설.
\item
  \textbf{{[}증명 스케치→5회차{]} Leech/\(\mathbb Q(\sqrt{13})\).} 위수 13 자기동형 \(g\), \(\eta=\sum(j/13)g^j\) (\(\eta^2=13\)). 최소벡터 \(w\) 의 \(s(w)=\tfrac12w\cdot\eta w\) 는 순환 Gram 행렬의 PSD 조건만으로 \(s\in\{0,\pm2,\pm4\}\) (\(\pm6\) 불가능: 직접 열거로 확인), 11-design 모멘트로 개수 \(a(0)=78624\), \(a(\pm2)=45864\), \(a(\pm4)=13104\). Hermite 노름 \(N(\nu(w))=13-s^2/4\in\{13,12,9\}\) 이고 단위원 논법으로 모든 \(w\ne0\) 에 \(N(\nu)\ge9=r^2\). ⇒ Leech 는 \(\Delta_{13}\) 위에서 \textbf{정확히 경계 안정}(\(R_U=1\))일 것으로 예상.
\end{itemize}

\hypertarget{uxd68cuxcc28-uxb9c8uxbb34uxb9ac-0741-utc-uxacbduxacfc-77uxbd84-uxbc18uxbc15-uxc5d0uxc774uxc804uxd2b8-3uxba85uxc758-uxd310uxc815}{%
\paragraph{3회차 마무리 (07:41 UTC, 경과 77분): 반박 에이전트 3명의 판정}\label{uxd68cuxcc28-uxb9c8uxbb34uxb9ac-0741-utc-uxacbduxacfc-77uxbd84-uxbc18uxbc15-uxc5d0uxc774uxc804uxd2b8-3uxba85uxc758-uxd310uxc815}}

\begin{itemize}
\tightlist
\item
  \textbf{확인됨}: 이득 항등식, 불변형식 차원 공식, \(q(x)\) 공식, 낮은 랭크 항을 포함한 평균값 공식과 \(T_U\) 점근식, 직선/판별식 기준, 엔트로피 부등식(창: \(q=3\): \(a/b\in[0.708,1.284]\) 등), Heisenberg \(p\le7\) 배제, \(SL_2(5)\times C_7\) 반례, \(g\le([G:C]/\chi(1))\,c/\varphi(c)\), Lemma N, Corollary N\textquotesingle, Prop U, 구면 사영 항등식.
\item
  \textbf{{[}반증/정정{]}}

  \begin{enumerate}
  \def\labelenumi{\arabic{enumi}.}
  \tightlist
  \item
    "\(\Delta=\mathbb Q\) 이면 군 자체가 배제된다"는 \textbf{텐서형 격자 \(W_\mathbb Z\otimes L\) 에만} 성립. 비극대 order 위의 불변 격자(예: \(D_8\) 에 대한 \(\{(x,y)\in L_1^2:y\equiv\phi(x)\bmod 2L_1\}\))는 모든 직선에서 경계 이상. 단 이는 자유 \(\mu_4/\mu_6\) 족의 부분족이라 이득은 같음. 모든 불변 격자에 유효한 강건한 no-go 는 소수 위수 원소의 엔트로피 기준.
  \item
    "희소 벡터 때문에 \(k=2\) 에서 방법이 깨진다"는 \textbf{인공물}: 반사원리 상계를 감금 반지름 밖까지 합산한 탓. 정지시간을 쓰면 모든 \(k\) 에서 무해. 또한 \textbf{대각합 없는(trace-free) 잡음}을 쓰면 \(V^{1+1/k}\) 손실 자체가 사라질 것으로 보임 {[}조건부{]} ⇒ Venkatesh 의 랭크 2 가군에서도 \(n^2m/\varphi(m)\) 이 가능할 수 있음.
  \item
    경계 항은 \(V\sim n^a\) 에서 \(T_U\sim n^{af-(1-f)/2}\) 로 \textbf{다항식적으로 클 수 있음} (\(V=n^2\) 에서 \(O(1)\) 이려면 \(f\le1/5\)).
  \item
    "모형 R(상수 1/2)"은 \(n\le10\) 자료로는 지지되지 않음: 실제 격자에서는 접촉이 일반 위치에 있지 않다(\(n=9,10\) 에서 정확히 \(N\) 쌍으로 동결되는 경우 0\%). ⇒ {[}휴리스틱{]}으로 강등. Poisson 모형의 "sticky 장벽"도 격자로 이전되지 않음(격자에서는 \(c_T\le\) 키싱수/2 뿐).
  \item
    GGS 접착 데이터의 "두 꼭대기" 예는 불가(소수 절단 보조정리: 각 소수 \(p\mid M\) 은 \(\varphi\)-질량의 \(>60\%\) 를 나눠야 함). "균형 잡힌 부채꼴" 은 미배제 → 7회차.
  \end{enumerate}
\item
  \textbf{새로 증명된 것(심사 중 추가)}: 접촉벡터가 \(\mathbb R^n\) 을 생성하면 단조 팽창 방향은 없다(\(B\le0\), \(Bx_i\cdot x_i=0\Rightarrow B=0\)). Hurwitz 가군에서는 \(q(x)\equiv1/4\) (상수) ⇒ Klartag 증명이 거의 그대로 적용될 것 → 6회차에서 완전 증명 시도.
\end{itemize}

\begin{center}\rule{0.5\linewidth}{0.5pt}\end{center}

\hypertarget{uxd68cuxcc28-07230757-utc-uxacbduxacfc-93uxbd84-uxd37cuxc149uxd2b8uxb860-uxbaa8uxd615-uxc774uxb860uxacfc-uxc218uxce58}{%
\subsubsection*{4회차 (07:23--07:57 UTC, 경과 93분): 퍼셉트론 모형 --- 이론과 수치}\label{uxd68cuxcc28-07230757-utc-uxacbduxacfc-93uxbd84-uxd37cuxc149uxd2b8uxb860-uxbaa8uxd615-uxc774uxb860uxacfc-uxc218uxce58}}

\textbf{선택한 과제}: 3회차 중 떠올린 "Klartag 과정 = 랜덤 선형계획/음의 퍼셉트론" 모형의 검증. (노트: \texttt{rounds/R4\_perceptron\_theory/notes.md}, \texttt{rounds/R4\_perceptron\_numerics/})

\textbf{{[}모형에서 증명{]} / {[}수치관찰{]}}
- 모형 재유도(상수 포함) 확인: 제약 \(\langle\theta,g_i\rangle\ge-\sigma_i\), 강도 \(\frac v4e^{\sigma/2}\), \(\mathrm{Var}(u^TB_0u)=2b^2/(n(n+2))\), \(V=ve^{\rho^2/8}\).
- \textbf{모형 정리 A (Gordon/CGMT)}: \(K_v\neq\emptyset\iff v<(2/e)N\); 임계에서 \(\rho^2=8\), 활성 제약 \(N\), \(V=2N=n(n+1)\). 구면 위에서는 단측: \(V\le\rho^2N/4\).
- 수치: 가우스 모형 LP \(N=100\): \(v^*/N=0.733\pm0.023\) (예측 0.7358), \(\rho^2=8.26\pm0.23\), \(V/N=2.03\pm0.03\). \textbf{실제 Poisson 점집합 (\(\mathbb R^n\), \(n=10..22\))의 대각합-LP: \(V_{LP}/N=1.98,1.99,1.95,2.09\)} (예측 2).
- \textbf{세 상수}: 증명된 \(1/(2e)\) (비사영 시계) → Klartag 과정 \(1/2\) (사영된 시계; 휴리스틱) → LP 꼭짓점 \(1\).
- \textbf{비볼록 영역}: 쌍 공식 \(P(\theta,\theta'\in K_v)=\exp(-V\Phi(\delta/4))\) {[}모형에서 증명{]}; Montanari--Zhong--Zhou(arXiv:2110.15824): 존재 임계/LP 임계 \(=\log|\kappa|\), 우리 변수로 \(\kappa=-\rho/2\) ⇒ \textbf{\(V_{exist}(\rho)=(2+o(1))N\log\rho\)} {[}조건부: 2차 모멘트 스케치{]}. \(\rho^2\sim16\log N\) 에서 \(\approx N\log\log N\), \(\rho=n^a\) 에서 \(\Theta(n^2\log n)\).
- 정정: 단순 2차 모멘트는 \(vV\ll\rho^2\) 일 때만 작동; "거리 \(O(1)\) 에서 탈상관"은 약한 의미로만 참.
- 모형의 한계: 정확한 Poisson 모형에서는 \(\sup V=\infty\)(바늘) ⇒ 절단 필요.

\textbf{막힌 지점(정확한 명제)}: (LP-1) 사영된 시계를 격자점에 대해 정당화하는 보조정리; (EX-1) 부피 \(\sim n^2\log n\) 타원체 쌍에 대한 구멍확률의 준-독립성; (EX-2) \(P_V\le e^{-cV}\) --- 이것 하나가 Venkatesh 추측을 증명함(Corollary N\textquotesingle). 현재 기술(Poisson 근사는 부피 \(\le cn\))로는 셋 다 도달 불가.

\hypertarget{uxd68cuxcc28-uxb9c8uxbb34uxb9ac-0757-utc-uxacbduxacfc-93uxbd84-uxc2ecuxc0ac-2uxba85uxc758-uxd310uxc815-uxc0c1uxc218uxb294-uxc804uxbd80-uxd655uxc778-uxd574uxc11duxc740-uxb300uxd3ed-uxc218uxc815}{%
\paragraph{4회차 마무리 (07:57 UTC, 경과 93분): 심사 2명의 판정 --- 상수는 전부 확인, 해석은 대폭 수정}\label{uxd68cuxcc28-uxb9c8uxbb34uxb9ac-0757-utc-uxacbduxacfc-93uxbd84-uxc2ecuxc0ac-2uxba85uxc758-uxd310uxc815-uxc0c1uxc218uxb294-uxc804uxbd80-uxd655uxc778-uxd574uxc11duxc740-uxb300uxd3ed-uxc218uxc815}}

\begin{itemize}
\tightlist
\item
  \textbf{확인됨(독립 재유도·독립 코드)}: 강도 \(\frac v4e^{\sigma/2}\), \(\mathrm{Var}(u^TB_0u)=2\|B_0\|^2/(n(n+2))\), AO 항등식 \(S(r)=4ve^{r^2/8}\)(분산 \(192V\)), Gordon 균형 \(4ve^{r^2/8}=r^2N\), \(v^*=(2/e)N\), 쌍 공식, MZZ 점근식(논문 본문 대조), Klartag sticky 과정 시뮬레이터의 정확성.
\item
  \textbf{{[}모형에서 증명됨, 신규 --- 심사자 발견{]} 정확한 \(\chi^2\) 법칙}: 정확한 Poisson 모형(모든 \(n\))과 가우스 모형(모든 \(N\))에서 대각합-LP 최적 타원체의 부피는 \textbf{정확히 \(V_{LP}\sim\chi^2_{2N}\)} (\(V/2\sim\mathrm{Gamma}(N,1)\); \(\mathbb EV=2N\), 표준편차 \(2\sqrt N\)). 증명 5줄: KKT ⇒ 유일한 \(N\)-부분집합이 결정하는 스케일 공변 정지집합, Mecke 공식. (KS 검정 \(p=0.31\)--\(0.65\).) ⇒ "\(V_{LP}\approx2N\)" 은 Gordon 도 가우스 근사도 필요 없는 정확한 정리이고, 표의 산포는 잡음이 아니라 이 법칙이다. Gordon 은 \(V=v^*\cdot e^{\rho^2/8}\) 의 분해(\(v^*=(2/e)N\), \(\rho^2\to8\))에만 필요.
\item
  \textbf{{[}증명됨{]} 격자에서의 대각합-LP = 쌍대 형식의 Voronoi 환원}: \(L=g\mathbb Z^n\) 의 대각합-LP 최적점은 쌍대 Gram 행렬 \((g^Tg)^{-1}\) 을 Voronoi 영역에 포함하는 \textbf{perfect form}. 따라서 \(V_{LP}(L)\le V_n\) 이고, \(n\le8\) 에서는 \(V_n<2N\) 이므로 Poisson 예측은 저차원 격자에 성립할 수 없다(Poisson LP 최적값은 \(n=4\) 에서 96\%, \(n=6\) 에서 100\% 가 실제 격자 최댓값을 초과).
\item
  \textbf{{[}증명됨{]}} 고정된 격자에서는 \(\sup_{\text{빈 타원체}}V=V_n<\infty\), Poisson 과정에서는 a.s. \(\infty\)(바늘). ⇒ 큰 이심률 영역(C4, \(n^2\log n\) 예측)의 Poisson 모형 진술은 격자로 직접 이전되지 않는다.
\item
  \textbf{{[}정정{]}} (i) 활성 제약의 offset 분포는 \(N(4,2.83^2)\) 가 아니라 \(r^2/2+rZ-\mathrm{Exp}(2)\) (평균 2, 표준편차 3.46). (ii) \(V=ve^{\rho^2/8}\) 는 2차 근사일 뿐: \(n\le26\) 에서 \(\log(V/v)=1.07\)--\(1.23\), \(v_{LP}/N=0.62\)--\(0.69\) --- "\(n=10\) 에서 이미 수렴" 은 두 15\% 오차의 상쇄. (iii) 임계 \(\|B_0\|^2=4(n+2)/n\). (iv) 절단 \(\sigma\le11\) 은 정확하지 않다(위반 기대 수 \(0.002N\)). (v) "가우스 모형은 \(\rho\ll n^{1/4}\) 에서 유효" 는 불충분(작용소 노름 조건 필요). (vi) \textbf{"알고리즘적 방법은 \(\Theta(N)\) 을 못 넘는다" 는 근거 없음}(MZZ 는 한 LP 알고리즘에 대해서만 증명) ⇒ 추측에서 \textbf{사변}으로 강등. (vii) "\(V_n\ge cn^2\log n\)" 은 절단 의존적 휴리스틱이지 모형이 지지하는 추측이 아니다. C4 상계 \(V\le CN\log\rho\) 도 그물 논법으로는 증명 안 됨(\(N\log(CN\rho^2)\) 뿐).
\item
  \textbf{심사자가 제시한 1차 모멘트 재정식화}: Poisson 에서 \(P(V_{LP}\le w)\le\mathbb E\#\{\)인증된 \(N\)-쌍 부분집합, 부피 \(\le w\}=(w/2)^N/N!\) 이고 이는 \(w<(2/e)N\) 에서 \(<1\) --- 즉 \textbf{모형에서는 순수 1차 모멘트만으로 Rezvyakova 상수의 2배(\(n^2/e\))}. 격자에서 빠진 보조정리: Haar-랜덤 격자에서 "공통 타원체 위의 eutaxy-인증 \(N\)-벡터쌍" 의 기대 개수 상계 (\(N\gg n\) 이라 Rogers 공식 밖).
\item
  \textbf{다음 할 일 → 9회차}: 실제 랜덤 격자에서 대각합-LP 를 계산해 \(V_{LP}(L)/N\) 을 측정(두 심사자 모두 "결정적 시험" 으로 지목).
\end{itemize}

\begin{center}\rule{0.5\linewidth}{0.5pt}\end{center}

\hypertarget{uxd68cuxcc28-07230759-utc-uxacbduxacfc-95uxbd84-sl_2p-uxd150uxc11cuxc871-leechmathbb-qsqrt13-uxc758-uxc815uxd655uxd55c-uxacc4uxc0b0}{%
\subsubsection*{\texorpdfstring{5회차 (07:23--07:59 UTC, 경과 95분): \(SL_2(p)\) 텐서족 --- Leech/\(\mathbb Q(\sqrt{13})\) 의 정확한 계산}{5회차 (07:23--07:59 UTC, 경과 95분): SL\_2(p) 텐서족 --- Leech/\textbackslash mathbb Q(\textbackslash sqrt\{13\}) 의 정확한 계산}}\label{uxd68cuxcc28-07230759-utc-uxacbduxacfc-95uxbd84-sl_2p-uxd150uxc11cuxc871-leechmathbb-qsqrt13-uxc758-uxc815uxd655uxd55c-uxacc4uxc0b0}}

\textbf{{[}증명됨(계산기, 정수 연산){]}} (\texttt{rounds/R5\_leech13/notes.md})
- 확장 Golay 부호로 Leech 격자 구성, \(M_{24}\) 생성원+부호변환+Conway \(\xi_T\) 로 위수 13 자기동형 \(g\) 획득 (고정점 없음, \(\mathrm{tr}=-2\)).
- \textbf{최소벡터 196560개 전수 분류}: \(s=-4,-2,0,2,4\) 에 대해 \(13104,45864,78624,45864,13104\); \(s=\pm6\) 은 0개 --- 0회차 예측과 정확히 일치.
- \(C(g)=\langle g\rangle\times SL_2(3)\) (위수 312); Leech 는 \(\Delta=\mathbb H\otimes\mathbb Q(\sqrt{13})\) 의 \textbf{극대 order} 위의 랭크 3 가군; \(\Delta\)-선형 자기동형군의 위수는 정확히 \textbf{2184}, 원소 위수 분포가 \(SL_2(13)\) 과 일치; 위수 3 원소는 대각합 0, 고정격자 \(\sqrt3E_8\) (\(\det=3^8\)).
- \textbf{정리(Leech 는 \(\Delta_{13}\) 위에서 정확히 경계 안정)}: 모든 \(\Delta\)-직선 \(U\) 에 대해 \(\det(\Lambda_{24}\cap U)\ge3^8\) (\(R_U\le1\)), 등호는 정확히 1638개 직선; 모든 \(\Delta\)-평면에 대해 \(R_U\le0.1001\).
- 이득: \(\sum_g\rho(g)^2=2184/72=30.33\), \(g=15.17\) --- icosian(7.5)의 2배. 이득의 대부분은 위수 7, 14 원소(비분리 토러스; 고정점 없음) 936개에서 나옴.

\textbf{본인 추가 유도 {[}증명됨{]}}: unimodular \(G\)-격자에서 고정점을 갖는 소수 위수 원소는 "정칙표현의 배수 + \(\mathbb Z[C_q]\)-자유"일 때만 허용되고 그때 정확히 \(R_U=1\) (Unimodular 강직성 보조정리). ⇒ \(SL_2(p)\) Weil 격자는 \((p+1)/2\) 가 소수일 때만 소수 위수 원소 시험을 통과 (\(p=5,13,37,61,73,157,\dots\); \(p=17,29,41\) 은 비분리 토러스의 위수 3 원소 때문에 실패).

\textbf{다음 할 일}: 5회차 이론 에이전트/심사 결과 반영; 일반 \(p\) 에서 \(N(\nu(w))\ge((p-1)/4)^2\) (직선 조건)이 \(\min(W_p)\ge\sqrt p\) 급 성장을 요구함을 정리.

\hypertarget{uxd68cuxcc28-uxb9c8uxbb34uxb9ac-0759-utc-uxacbduxacfc-95uxbd84-uxc774uxb860-uxc5d0uxc774uxc804uxd2b8-uxc2ecuxc0ac-2uxba85}{%
\paragraph{5회차 마무리 (07:59 UTC, 경과 95분): 이론 에이전트 + 심사 2명}\label{uxd68cuxcc28-uxb9c8uxbb34uxb9ac-0759-utc-uxacbduxacfc-95uxbd84-uxc774uxb860-uxc5d0uxc774uxc804uxd2b8-uxc2ecuxc0ac-2uxba85}}

\begin{itemize}
\tightlist
\item
  \textbf{심사 판정}: R5 의 모든 주장이 독립 코드(자체 Golay 부호, 자체 HNF, 정수 연산)로 재현됨. "Leech 는 \(\Delta_{13}\) 위에서 정확히 경계 안정" 은 \textbf{계산기 보조 정리로 확정}. 추가로 (a) 직선마다 최소벡터 24개(630개 궤도 대표 전수), (b) \textbf{경계 직선이 정확히 \(1638=3\times546\) 개}임을 class number 없이 증명(Hermite 상수 \(\gamma_8=2\), \(E_8\) 의 유일성, 7-design): 세 종류 = \(s=+4\) 최소벡터 직선 546, \(s=-4\) 직선 546, \(\sqrt3E_8\) 형 546(그중 91개가 위수 3 원소의 고정격자), (c) 노름 6·8 벡터 전체의 \(N(\nu)\) 분포(최소 9, 16), (d) class number 1 (\(O^1=SL_2(3)\) 정확히; Eichler mass \(1/12\)).
\item
  \textbf{일반 \(p\) {[}증명됨{]}} (\texttt{rounds/R5\_sl2p\_theory/notes.md}): Lemma A′(unimodular 강직성; 본인 유도와 독립적으로 일치), Cor. A″(\((p+1)/2\) 합성수 ⇒ 불안정: \(R_U\ge9.99\) (\(p=17\)), \(11.8\) (\(29\)), \(12.5\) (\(41\))), Lemma B(분리 토러스의 홀수 위수 부분군의 고정격자는 정확히 \(R=1\)), \((p+1)/2\) 소수일 때 고정공간 분류, 직선 기준 \(N(\nu(w))\ge r^2\). \textbf{심사자 강화}: unimodular 가 아니어도 \(R_U(L)R_U(L^\#)=((m/a)^aq^{-c})^2\ge1\) --- 즉 쌍대가 같은 족에 있는 모든 \(G\)-격자에 Cor. A″ 성립.
\item
  \textbf{{[}정정{]}} (i) \(\Delta=(-1,-1)_F\) 는 \(p\equiv5\pmod8\) 에서만 참(\(p\equiv1\pmod 8\) 에서는 2 위에서 분기하는 다른 대수). (ii) "최소노름만으로는 결코 충분치 않다" 는 과장(특정 두 부등식이 부족하다는 것만 증명). (iii) "텐서 불안정" 은 \textbf{방법(1차 모멘트/Klartag 형 논증)의 장애}이지 그 족에 조밀한 격자가 없다는 증명이 아니다(결정론적 모순은 \(R_U>4^{\dim U}\) 필요). (iv) "\(\approx15\times\) Klartag (\(2.79n^2\), \(n=24k\))" 은 \textbf{정리가 아니라 휴리스틱}: 1638개 경계 직선 각각이 \(V=cn^2\) 에서 \(\approx1.18c^{1/3}n^{1/3}\) 개의 격자점을 공 안에 가지며, 순수 1차 모멘트 논증은 이 족에서 \(V\sim n\) 에서 멈춘다. Klartag 형 과정에서는 이 점들이 접촉 궤도 \(\sim n^{1/3}\) 개/직선을 추가할 뿐(\(D_k\sim n^2/144\) 에 비해 무시 가능)일 것으로 보이나 {[}휴리스틱{]} 증명 없음.
\item
  \textbf{미해결로 남은 것}: \(p=37,61,73,157,\dots\) (\((p+1)/2\) 소수). 일반적(generic) 격자에 대한 Siegel 형 기대 위반 벡터 수(단위군 궤도당): \(1.6\times10^5\) (\(p=37\)), \(246\) (\(61\)), \(2.3\) (\(73\)), \(2\times10^{-26}\) (\(157\)). → \textbf{10회차에서 \(p=37,61\) 을 직접 계산}.
\end{itemize}

\begin{center}\rule{0.5\linewidth}{0.5pt}\end{center}

\hypertarget{uxd68cuxcc28-07310808-utc-uxacbduxacfc-104uxbd84-hurwitzicosian-uxac00uxad70-uxaca9uxc790uxc5d0-uxb300uxd55c-uxba85uxc2dcuxc801-uxc0c1uxc218uxc758-klartag-uxd615-uxc815uxb9ac-uxc644uxc804-uxc99duxba85-uxc2ecuxc0ac-2-uxac1cuxc815}{%
\subsubsection*{6회차 (07:31--08:08 UTC, 경과 104분): Hurwitz/icosian 가군 격자에 대한 명시적 상수의 Klartag 형 정리 --- 완전 증명 + 심사 2 + 개정}\label{uxd68cuxcc28-07310808-utc-uxacbduxacfc-104uxbd84-hurwitzicosian-uxac00uxad70-uxaca9uxc790uxc5d0-uxb300uxd55c-uxba85uxc2dcuxc801-uxc0c1uxc218uxc758-klartag-uxd615-uxc815uxb9ac-uxc644uxc804-uxc99duxba85-uxc2ecuxc0ac-2-uxac1cuxc815}}

\textbf{선택한 과제}: 3회차 심사에서 나온 관찰(\(q(x)\equiv1/4\))을 이용해 자유 작용 대칭족에서 상수까지 추적한 엄밀한 정리를 쓴다. (\texttt{rounds/R6\_quaternionic/proof\_revised.md}, 317줄)

\textbf{{[}증명됨 --- 인용 입력에 조건부{]}} (입력: (L1) Gargava arXiv:2107.04844 Thm 16 의 평균값 정리, (L2) Klartag 의 보조정리들, (L3) 표준 확률해석)
- \textbf{정리 C′ (틀)}: \(D\) 가 \(\mathbb Q\)-나눗셈대수, \(D_\mathbb R\cong K^s\) (\(K\in\{\mathbb R,\mathbb C,\mathbb H\}\)), \(G\subset D^\times\) 유한군으로 \(\mathbb Q\)-span 이 \(D\), \(w=|G|\), \(\delta=\dim_\mathbb QD\) 이면 \(k\to\infty\) 에서
\[V_{k\delta}\ \ge\ (1-o(1))\,\frac{w}{4e\,\delta}\,n^2\;=\;(1-o(1))\,\frac{w\,D_k}{2e}.\]
- \textbf{정리 A (Hurwitz, \(4\mid n\))}: \(V_n\ge(\tfrac3{2e}-o(1))n^2\approx0.552\,n^2\). \textbf{정리 B (icosian, \(8\mid n\))}: \(V_n\ge(\tfrac{15}{4e}-o(1))n^2\approx1.380\,n^2\). (자명군: Rezvyakova 의 \(\tfrac1{2e}=0.184\) 재현. Eisenstein, \(2\mid n\): \(\tfrac3{4e}\).)
- 핵심 단계: (1) 불변 형식 = Hermite 행렬 \(L_H\), \(P(x\otimes x)=L_{xx^*/d}\), \(q(x)=\frac1d\sum_vu_v^2\ge q_0=1/\delta\) (Hurwitz 에서는 상수 \(1/4\)); (2) \textbf{주변(ambient) Dyson 운동을 유지}하고 제약공간만 \(F_A\subset\mathcal Q\) 로 줄이면 Klartag §2, Lemma 3.1/3.3, Prop 4.2 가 그대로 적용; (3) 자유 궤도 ⇒ \(\dim F_t\ge D_k-|\text{접촉}|/w\); (4) Rezvyakova 적분 보조정리 + \textbf{각도 보조정리}(Dirichlet 분포; icosian 의 방향 의존성 처리); (5) 최적화 \(e^z=wq_0n^2e^\rho/4\).
- \textbf{심사 2명}: "수학적 오류 없음, 상수 \(3/(2e)\), \(15/(4e)\) 확인(독립 재유도 + 구적: \(n=10^4,10^5,10^6\) 에서 극한값 대비 1.19, 1.045, 1.02 배), 사소한 수정만 필요". Klartag 의 {[}V{]} 단계들을 원문과 한 줄씩 대조.
- \textbf{한계(정직하게)}: (i) \(o(1)\) 은 비효과적이고 느리다(\((\log n)^{-1/4}\)). (ii) \(4\mid n\), \(8\mid n\) 부분열에서만. (iii) \textbf{차수는 \(n^2\) 그대로} --- \(n^2\log\log n\) (AGZ/GGS) 보다 점근적으로 약하다. 다만 원분체 이득 \(m/(2\varphi(m))\) 이 \(7.5\) 를 넘으려면 \(\log\log m>8\) 이므로, 현실적인 모든 차원에서 icosian 상수가 더 크다. (iv) Gargava Thm 16 의 증명 내부는 재검증하지 않음(개정자가 \(D_\mathbb R=K^s\), \(s\ge2\) 에서 "추이성" 표현의 문구상 문제를 지적: 열린 조밀 궤도로 대체하면 될 것으로 보이나 미검증). (v) 새로움: AGZ/GGS 와의 관계(그들의 틀이 비가환 \(D\) 를 이미 포함하는지)는 확인하지 못함.
- 개정자가 일반 행렬 블록(\(D_\mathbb R\supset M_m(K)\)) 경우를 위한 §7(Wishart 형 각도 보조정리)을 추가했으나 \textbf{미심사} ⇒ {[}조건부{]}.

\begin{center}\rule{0.5\linewidth}{0.5pt}\end{center}

\hypertarget{uxd68cuxcc28-07320800-utc-uxacbduxacfc-96uxbd84-ggs-uxc6d0uxbd84-uxc811uxcc29-uxd2c0-uxc548uxc5d0uxc11c-mcmvarphim-uxc774-uxac00uxb2a5uxd55cuxac00}{%
\subsubsection*{\texorpdfstring{7회차 (07:32--08:00 UTC, 경과 96분): GGS 원분 접착 틀 안에서 \(M/C>M/\varphi(M)\) 이 가능한가}{7회차 (07:32--08:00 UTC, 경과 96분): GGS 원분 접착 틀 안에서 M/C\textgreater M/\textbackslash varphi(M) 이 가능한가}}\label{uxd68cuxcc28-07320800-utc-uxacbduxacfc-96uxbd84-ggs-uxc6d0uxbd84-uxc811uxcc29-uxd2c0-uxc548uxc5d0uxc11c-mcmvarphim-uxc774-uxac00uxb2a5uxd55cuxac00}}

\textbf{결과}: 개선 없음, 완전한 no-go 도 없음. (\texttt{rounds/R7\_gluing/notes.md}, \texttt{referee/})
- \textbf{{[}증명됨{]}} 재정식화: \(M/C=1/\mu_M(\mathcal N)\) (\(\mu_M(n)=\varphi(n)/M\) 은 약수 격자 위의 곱측도) \(=\sum_g\rho(g)^2\) (Parseval). admissibility 의 지지집합 부분은 정확히 LP(협력게임의 core); 절단·분할·꼭짓점 조건; \textbf{소수별 창} \(x_q<x_{hi}(q)<2/q\).
- \textbf{{[}증명됨{]} 정리 A (부분 no-go, 홀수 \(M\))}: \(M/C>M/\varphi(M)\) ⇒ \(M\notin\mathcal N\) 이고 \(\sum_{q\mid M}x_{hi}(q)>1\). 따라서 홀수 소인수 3개 이하, 또는 4개인데 \(\{3,5,7,11\}\) 이 아니면 불가능. \textbf{정리 B}: \(M/C\le\frac M{\varphi(M)}\prod_p(p-1)^{x_{hi}(p)}\) (약함).
- \textbf{{[}수치관찰{]}} 전수 LP 탐색: 약수 24개 이하의 모든 시험 모듈러스(홀수 13개 + 심사자의 짝수 34개)에서 \(M/C>M/\varphi(M)\) 인 admissible 집합 \textbf{0개}; lcm 을 피하는 집합은 소수별 절단조차 통과 못 함.
- \textbf{{[}심사 정정{]}} 정리 A 는 홀수 \(M\) 에만; "모든 점근족이 선행 차수에서 실패" 는 과장(큐브−꼭대기 족은 선행 차수에서 정확히 경계); "비순환 아벨군은 도움이 안 된다" 는 근거 없음(미해결로 표기).
- \textbf{{[}추측{]}} 모든 admissible 데이터에 \(M/C\le K\cdot M/\varphi(M)\). \textbf{막힌 지점}: admissibility 부등식은 모두 "엔트로피 ≤ 경계 가중치" 꼴이라 \(\mathcal N\) 이 가벼울수록 쉬워진다 --- 상계는 피복조건과 소수별 창의 결합에서만 나오며, 점근 영역(\(\sum1/q\sim\log\log y\))에서는 닫히지 않는다.

\begin{center}\rule{0.5\linewidth}{0.5pt}\end{center}

\hypertarget{uxd68cuxcc28-07380816-utc-uxacbduxacfc-112uxbd84-uxb4f1uxbcc0equivariant-klartagvoronoi-uxd0d0uxc0c9-uxb300uxce6d-uxacfcuxc815uxc740-uxc5b4uxb5a4-uxaca9uxc790uxb97c-uxcc3euxb294uxac00}{%
\subsubsection*{8회차 (07:38--08:16 UTC, 경과 112분): 등변(equivariant) Klartag/Voronoi 탐색 --- 대칭 과정은 어떤 격자를 찾는가}\label{uxd68cuxcc28-07380816-utc-uxacbduxacfc-112uxbd84-uxb4f1uxbcc0equivariant-klartagvoronoi-uxd0d0uxc0c9-uxb300uxce6d-uxacfcuxc815uxc740-uxc5b4uxb5a4-uxaca9uxc790uxb97c-uxcc3euxb294uxac00}}

(사용자 제안 "작은 \(n\) 에서 실례를 직접 찾고 규칙성을 본다" 의 대칭판. \texttt{rounds/R8\_equivariant/}, 1152회 실행, 검증자가 전부 재검증)
- \textbf{{[}수치관찰, 검증됨{]} 기록 격자들이 그대로 나온다}: \(E_8\) (Hurwitz \(k=2\), icosian \(k=1\), \(\mathbb Z[\zeta_5]\) \(k=2\): 100\%), \(K_{12}\) (\(\mathbb Z[\zeta_7]\) \(k=2\): 33\%), \(BW_{16}\) (icosian \(k=2\): 42--67\%, Hurwitz \(k=4\): 12--40\%), \(\Lambda_{20}\) (Hurwitz \(k=5\): 2/10), \textbf{Leech (icosian \(k=3\), 자유도 \(D=30\) 뿐: 11/24 실행)}. 검증자가 정수 연산으로 확인(짝 unimodular, 최소 4, 196560개 ⇒ Leech).
- \textbf{{[}수치관찰{]} 스케일링 법칙}: 첫 \(\mathcal Q\)-perfect 형식에서 \(V_{perf}\approx1.06\times(wD/2)\) (17개 (족,\(k\)) 점, \(wD/2=50..1800\), 로그-로그 기울기 0.94). 즉 비대칭의 "\(V\approx N\)" 이 \textbf{"\(V\approx\) (궤도 크기/2)×(불변형식 차원)"} 으로 일반화 --- 6회차 정리의 상수 구조 \(wD_k/(2e)\) 와 같은 형태(상수 \(1\) vs \(1/e\)).
- 각 차원에서 족의 순위는 \(wD/2\) 순(예: \(n=24\): icosian \textgreater{} Hurwitz \textgreater{} \(\mathbb Z[\zeta_{13}]\)).
- 발견된 모든 \(\mathcal Q\)-perfect 형식은 유리 형식(분모 \(q=2\): root 격자, \(q=4\): \(\Lambda_{10},K_{12},BW_{16},\Lambda_{20}\), Leech). 최선 형식의 정수 Gram 행렬: \texttt{best\_forms.json}.
- \textbf{한계}: 알려진 최선값을 넘는 \(V\) 는 없음(max \(V/V_{best}=1.000000\)). 소차원 기록 격자들은 극도로 퇴화된 꼭짓점(Leech: 1638 궤도 vs \(D=30\))이며 점근 영역(\(V\sim wD\))과는 다른 체제.
- 버그 1건(탐욕 LP 하강이 약한 eutactic 비극단 형식에서 멈춤; 7/1152) --- 보고된 수치를 보수적으로만 만든다.

\begin{center}\rule{0.5\linewidth}{0.5pt}\end{center}

\hypertarget{uxd68cuxcc28-07550821-utc-uxacbduxacfc-117uxbd84-uxc2e4uxc81c-uxb79cuxb364-uxaca9uxc790uxc5d0uxc11cuxc758-uxb300uxac01uxd569-lp-poisson-uxc608uxce21-v_lp2n-uxc758-uxc2dcuxd5d8}{%
\subsubsection*{\texorpdfstring{9회차 (07:55--08:21 UTC, 경과 117분): 실제 랜덤 격자에서의 대각합-LP --- Poisson 예측 \(V_{LP}=2N\) 의 시험}{9회차 (07:55--08:21 UTC, 경과 117분): 실제 랜덤 격자에서의 대각합-LP --- Poisson 예측 V\_\{LP\}=2N 의 시험}}\label{uxd68cuxcc28-07550821-utc-uxacbduxacfc-117uxbd84-uxc2e4uxc81c-uxb79cuxb364-uxaca9uxc790uxc5d0uxc11cuxc758-uxb300uxac01uxd569-lp-poisson-uxc608uxce21-v_lp2n-uxc758-uxc2dcuxd5d8}}

(\texttt{rounds/R9\_trace\_lp\_lattices/}; 432개 격자의 LP 를 쌍대 인증서로 검증)
- \textbf{{[}증명됨{]}} 대각합-LP 최적점 = \(Q^{-1}\) 을 Voronoi 영역에 갖는 perfect form (LP 쌍대성). 검증: \(A_6,D_6,E_6,E_7,E_8,D_4,A_{10}\) 에서 LP 가 자기 자신을 반환.
- \textbf{{[}수치관찰{]} 실제 격자는 \(n\le28\) 에서 Poisson 모형과 다르다.} (a) \(n=6,7,8\): LP 꼭짓점이 대부분 \(E_6\) (43/48), \(E_7\) (23/24), \(E_8\) (38/48). (b) \(n=9\)--\(13\): 꼭짓점이 단순(정확히 \(N\) 쌍)인 경우 0/172, \(V_{LP}/N=1.16\)--\(1.27\), \(v_{LP}/N=0.42\)--\(0.48\). (c) \(n\ge14\): 꼭짓점이 단순해지지만 \(V_{LP}/N\) 은 \(1.20\) (\(n=14\)), \(1.28\) (\(16\)), \(1.35\) (\(18\)), \(1.44\) (\(20\)), \(1.55\) (\(22\)--\(24\)).
- \textbf{{[}검증자 정정{]}} \(p=100003\) Goldstein--Mayer 족은 \(n\ge14\) 부터 입도(granularity) 문제가 있고 \(n=28\) 행(\(1.44\))은 인공물. 큰 소수(\(p=10^{12}+39\))로 다시 하면 \(V_{LP}/N=1.28\) (\(n=16\)), \(1.44\) (\(20\)), \(1.55\) (\(24\)), \textbf{\(1.75\pm0.05\) (\(n=28\))} --- 여전히 증가 중. ⇒ "\(V_{LP}(L)/N\to2\) (\(n\to\infty\))" 는 \textbf{기각되지 않았고 결정되지도 않았다} {[}추측{]}.
- 같은 격자에서 Klartag 형 sticky walk: \(V_K/N=1.10,1.26,1.02\)--\(1.05,1.06\) (\(n=6,8,10,12\)); LP 가 대체로 13--20\% 더 크다.
- \textbf{교훈}: 격자의 접촉벡터는 \(n\) 개 기저의 정수결합이고 kissing 배치를 이루므로 Mecke/Gamma 논법이 실패 --- 저차원에서는 root 격자의 뚱뚱한(비단체) Voronoi 영역이 지배한다. 상수 개선(\(1/(2e)\to\) LP 의 \(1\))에 필요한 보조정리(4회차 LP-1)는 여전히 미해결.
- \textbf{다음 할 일}: (10회차) 유일하게 남은 "차수 개선" 후보인 \(SL_2(p)\) 텐서족의 \(p=37,61\) 을 직접 계산으로 시험; (11회차) 그 후보의 Mordell--Weil 재정식화 검증과 조건부 주장에 대한 반박 시도.

\begin{center}\rule{0.5\linewidth}{0.5pt}\end{center}

\hypertarget{uxd68cuxcc28-08260846-utc-uxacbduxacfc-142uxbd84-sl_2p-uxd6c4uxbcf4uxc758-mordellweil-uxc7acuxc815uxc2dduxd654-uxac80uxc99dt1-uxc870uxac74uxbd80-uxc8fcuxc7a5uxc5d0-uxb300uxd55c-uxbc18uxbc15-uxc2dcuxb3c4t2}{%
\subsubsection*{\texorpdfstring{11회차 (08:26--08:46 UTC, 경과 142분): \(SL_2(p)\) 후보의 Mordell--Weil 재정식화 검증(T1) + 조건부 주장에 대한 반박 시도(T2)}{11회차 (08:26--08:46 UTC, 경과 142분): SL\_2(p) 후보의 Mordell--Weil 재정식화 검증(T1) + 조건부 주장에 대한 반박 시도(T2)}}\label{uxd68cuxcc28-08260846-utc-uxacbduxacfc-142uxbd84-sl_2p-uxd6c4uxbcf4uxc758-mordellweil-uxc7acuxc815uxc2dduxd654-uxac80uxc99dt1-uxc870uxac74uxbd80-uxc8fcuxc7a5uxc5d0-uxb300uxd55c-uxbc18uxbc15-uxc2dcuxb3c4t2}}

\textbf{선택한 과제}: 9회차까지 살아남은 유일한 차수 개선 후보를 (a) 대수기하의 언어로 재정식화하고 (b) "성립하면 \(n^{2.5}\)" 라는 조건부 주장 자체를 공격한다. (\texttt{rounds/R11\_mordell\_weil/T1,T2/notes.md})

\textbf{T1 --- Mordell--Weil 재정식화 {[}증명됨; 일부 문헌은 접근 불가라 기억에 의존{]}} (출발점은 본인 유도, T1 이 한 줄씩 검증·정정)
- (M1) \(C_p:\ y^2=x^p-x\) (\(g=(p-1)/2\)) 위에 \(SL_2(p)\) 가 \(x\mapsto\frac{ax+b}{cx+d}\), \(y\mapsto y/(cx+d)^{(p+1)/2}\) 로 충실하게 작용(\(-1\) = 초타원 대합), \(H^1=\psi+\psi'\) (위수 \(p\) 원소의 Lefschetz 수 \(2-3=-1\); 비분리 토러스 생성원의 대각합 2 로 이산계열 표현 배제).
- (M2) \(\mathrm{Jac}(C_p)\) 는 superspecial(Cartier 작용소가 \(x^idx/y\) 를 죽임), Frobenius \(\pi\) 는 \(\pi^2=p\) (\(\#C_p(\mathbb F_p)=\#C_p(\mathbb F_{p^2})=p+1\)) 이고 \(G\) 와 가환 ⇒ \textbf{\(\eta=\pm\pi\): \(F=\mathbb Q(\sqrt p)\)-구조는 Frobenius 가 준다.} \(M_p:=\mathrm{Hom}(\mathrm{Jac}\,C_p,E)\) 는 랭크 \(2(p-1)\) 의 \(G\)-격자이고 \(\mathrm{End}_G(M_p)\) 는 극대 order.
- (M3) \textbf{짝 unimodular Weil 격자 \(W_p\) (\(E_8\), Leech, \ldots) 는 \(M_p\) 에 뒤틀린 형식 \(\mathrm{Tr}_{F/\mathbb Q}(\tfrac\varepsilon{\sqrt p}\beta)\) 를 준 것}; \(G\)-불변 형식의 원 \(Q_t=e^tm_++e^{-t}m_-\) (주기 \(2\log\varepsilon\)) 위에서 unimodular 점과 Mordell--Weil 점(\(2\deg=\sqrt p\,Q\))은 정반대(거리 \(\log\varepsilon\)).
- (M4) \textbf{직선 안정성 \(\iff N_{F/\mathbb Q}(q_F(\phi))\ge r^2=((p-1)/4)^2\) (\(\forall\phi\ne0\))}, \(q_F\) 는 \(\mathrm{Tr}\,q_F=\deg\) 인 \(O_F\)-값 이차형식. 특히 모든 비상수 사상 \(C_p\to E\) 의 차수가 \(\ge(p-1)/2=g\) 여야 한다. 초등적 한계(Castelnuovo--Severi)는 \(\deg\ge(p+3)/4\) 뿐. 랭크 \(j\): \(N(\mathrm{Mdet}\,H_U)\ge(r/j)^{2j}\).
- (M5) \(p\equiv5\pmod 8\): \(x\mapsto x^{(p-1)/4}\) 로 \(C_p\to D:Y^2=X^5-X\) (차수 \((p-1)/4\)), \(D\to E\) (차수 2) ⇒ 차수 정확히 \((p-1)/2\) 인 사상 = 경계 벡터. {[}정정{]} \(3\mid p+1\) 이면 차수 \((p+1)/3<(p-1)/2\) 인 Galois 타원 몫이 있다(\(p=17,29,41\) --- 모두 \((p+1)/2\) 합성수, 5회차와 일치). (M6) 차수 \(d\) 사상 ↔ \(R^*(e_1+\dots+e_4)\equiv\sum_{P\in\mathbb P^1(\mathbb F_p)}P\pmod2\) 인 차수 \(d\) 유리함수 \(R\).
- (M7) {[}휴리스틱{]} generic 격자의 기대 위반 벡터 수(단위군 궤도당) \(=h\,\omega_h^2c^{h/2}\log\varepsilon\): \(1.15\times10^4\) (\(p=13\)), \(1.58\times10^5\) (\(37\)), \(246\) (\(61\)), \(2.3\) (\(73\)), \(1.8\times10^{-26}\) (\(157\)); \(G\)-궤도 수로는 \(5.3,\ 3.1,\ 10^{-3},\ 6\times10^{-6},\dots\) ⇒ \textbf{\(p=37\) 은 불안정이 예상되고 \(p\ge61\) 은 구조적 짧은 벡터가 없는 한 안정이 예상됨.} (단 \(p=13\) 에서 이미 과대예측: 실제 위반 0.) 랭크 2 이상의 부분공간은 휴리스틱상 무관.
- \textbf{Venkatesh 의 추측}의 출처: AGZ (arXiv:2606.05105) 서론에 "On various occasions, Venkatesh has conjectured \(\Delta_N^{lat}\le N^2(\log N)^{O(1)}2^{-N}\)" 로 인용(Venkatesh 본인의 문헌 인용은 없음).
- 문헌 한계: Gross (JAMS 1990), Nebe (J. Algebra 1998), Dummigan, Watkins 는 네트워크 정책상 열람 불가 --- "Gross 가 이 곡선과 \(p=13\)⇒Leech 를 명시했는지", \(p=17,29,37,61\) 격자의 최소노름 문헌값은 확인 못 함.

\textbf{T2 --- 반박 시도 {[}치명적 장애 없음, 그러나 주장은 수정 필요{]}}
- \textbf{{[}증명됨{]} 텐서족에서 전형적 사영 분산은 \(q=q_0(1+\tfrac{r-1/2}{k})\)} (Monte Carlo 확인; 초과분 \(r/k=n/(2D_k)\) 는 정확히 항등 방향(호흡 모드)의 기여). ⇒ 기존 회계로는 \(V_{final}=V_0^{r/(k+r)}(|G|D_k/2e)^{k/(k+r)}\): \textbf{\(k\approx r\) 에서는 이득이 없다(\(n^{2.0}\))}; \(k=cr\) 이면 지수 \(2+\tfrac{c-1}{2(c+1)}\); 완전한 이득 \(p\) 는 \(k\ge Cr\log n\) 에서만 ⇒ \(n\sim Cp^2\log p\), \(V\sim n^{2.5}/\sqrt{\log n}\) (여전히 \(n^2\mathrm{polylog}\) 와 충돌). \(k=r+1\) 판은 대각합 없는 잡음을 쓰는 새 변형(미증명)을 요구.
- (A1) 평균값 공식 주항 상수는 정확히 1; \(O\)-가군 제타 \(\prod_{i<2r}\zeta_F(2k-i)\) 는 \(k\ge r+1\) 에서 수렴(\(\le1.17\)). (A2) 강제되는 경계류(분리 토러스 고정공간, \(SL_2(3)\)-isotypic 성분)가 흡수하는 제약 궤도는 \(O(p)\) 로 \(D_k\sim p^2/4\) 에 비해 무시 가능; 단 경계 직선의 개수 상계는 추가 가정. (A5) 상계 쪽 장애는 발견 못 함(theta 급수의 Hilbert 모듈러 제약은 미검토).
- \textbf{{[}신규 필요조건{]}} 부분군 \(H\) 의 \(\mathbb Q\)-isotypic 성분도 \(\Delta\)-부분공간이다. \(SL_2(3)\) 의 Hurwitz 성분: \(R=2^{-e}\le1\) (항상 안전). \textbf{\(p\equiv1\pmod5\) (\(p=61,421,\dots\)) 에서 \(SL_2(5)\) 의 icosian 성분: \(R=(625/256)^{2r/15}2^{-e}\) --- \(p=61\) 은 2-진 접착 지수 \(e\ge3\) 이어야 안정} (미계산 → 12회차).
- \textbf{수정된 조건부 주장}: (H1) 무한히 많은 \(p\) 에서 모든 \(\Delta\)-부분공간에 \(R_U\le1\), (H2) 사실상 경계인 부분공간들의 접촉 기여가 \(o(D_k)\), (H3) \(r\ge2\) 텐서족에 대한 등변 Klartag 논증(\(k\ge Cr\log n\)) ⇒ \(n=8rk\sim\tfrac C2p^2\log p\) 에서 \(V_n\ge(1-o(1))\tfrac{p\,n^2}{2e}\ge c\,n^{5/2}/\sqrt{\log n}\).

\begin{center}\rule{0.5\linewidth}{0.5pt}\end{center}

\hypertarget{uxd68cuxcc28-0824uxc9c4uxd589-uxc911-sl_2p-weil-uxaca9uxc790uxc758-uxc9c1uxc811-uxad6cuxc131uxacfc-uxc548uxc815uxc131-uxc2dcuxd5d8-p51317293761}{%
\subsubsection*{\texorpdfstring{10회차 (08:24--진행 중): \(SL_2(p)\) Weil 격자의 직접 구성과 안정성 시험 (\(p=5,13,17,29,37,61\))}{10회차 (08:24--진행 중): SL\_2(p) Weil 격자의 직접 구성과 안정성 시험 (p=5,13,17,29,37,61)}}\label{uxd68cuxcc28-0824uxc9c4uxd589-uxc911-sl_2p-weil-uxaca9uxc790uxc758-uxc9c1uxc811-uxad6cuxc131uxacfc-uxc548uxc815uxc131-uxc2dcuxd5d8-p51317293761}}

\textbf{선택한 과제}: 5회차·11회차의 "첫 미해결 사례 \(p=37\)" 을 계산으로 결정한다. 독립 구현 3개: 빌더 A, 빌더 B(하위 에이전트), 그리고 본인(오케스트레이터)의 프로토타입 \texttt{rounds/R10\_weil\_lattices/orchestrator/} (\texttt{proto.py}, \texttt{run37.py}, \texttt{sub37.py}, \texttt{certify.py}).

\textbf{방법(본인 유도 → 세 구현이 검증)}: 홀함수 Weil 표현 \(T=\mathrm{diag}(\zeta^{x^2})\), \(S=\kappa(\zeta^{2xy}-\zeta^{-2xy})/\sqrt p\) (\(\kappa=-1\) 에서 \(S^2=(ST)^3=-I\))을 \(\mathbb Q(\zeta_p)\) 의 모든 매장으로 Minkowski 매장. \(v_0=\mathbf 1_{\mu_m}-\mathbf 1_{-\mu_m}\) (\(m=(p-1)/4\); 분리 토러스 \(C_m\) 의 고정벡터라서 \(\mathbb Q[G]v_0\cong W\) 한 벌). \(\Lambda=\mathbb Z[G]v_0\), \(B=\mathrm{Tr}_{K/\mathbb Q}\langle\cdot,\cdot\rangle\) (정수값), \(m_\pm\) = 매장을 \((a/p)=\pm1\) 로 나눈 노름. \(D=\det_B(\Lambda)^{1/(p-1)}\), \(\theta_p=(p-1)^2/(4p)\), \(\Gamma(w)=m_+(w)m_-(w)/D\). \textbf{\(\Gamma(w)<\theta_p\) ⇒ \(\Delta w\) 직선에서 \(R_U\ge(\theta_p/\Gamma)^4>1\).}

\textbf{결과}
- \textbf{검증} \(p=5\): \(D=5\), \(\min\Gamma=0.8=\theta_5\); unimodular 점(\(t=\pm\log\varepsilon\))이 \(E_8\). \(p=13\): \(D=117=r^2p\), \textbf{\(\min\Gamma=36/13=\theta_{13}\) 정확히}(전수 열거; 세 류 각각 6552쌍 --- 5회차의 13104+13104+26208/2 와 일치), unimodular 점이 Leech (최소 4). \(C_3\)-고정격자: \(\log R_U=0.000000\), \(\sqrt3E_8\) (노름 6: \(N(\nu)=9\) 가 24쌍, \(27\) 이 96쌍 --- 5회차와 일치). 직접 인증: \(O\cdot w\) 의 행렬식 \(=((p/4)m_+m_-)^4\), \(R_U=1.000000\).
- \textbf{\(p=17,29\) (예상대로 불안정)}: 빌더 B: \(p=29\) 비분리 위수 5 원소의 고정격자가 \(\Delta\)-직선이고 \(R_U=(49/25)^4=14.76\); 위수 3: \(R_U=11.78\).
- \textbf{{[}수치(부동소수점, 잔차 \(\le2\times10^{-9}\)), 독립 구현 2개 일치{]} \(p=37\) 은 텐서 불안정.} \(\Lambda\): 랭크 72, \(D=2997=81\cdot37\); unimodular 점은 72차원 짝 unimodular 격자 \(W_{37}\). 분리 토러스 \(C_9\) 의 고정 직선은 정확히 경계(\(\log R_U=0\), \(\min\Gamma=\theta\)) --- 보정 확인. 그러나 \textbf{\(C_3\)-고정 부분격자(랭크 24, 그 자체는 \(R_U=1\))는 \(\sqrt3\times\)(근 144개를 갖는 unimodular 랭크 24 격자, 즉 뿌리계 \(D_4^6\) 또는 \(A_5^4D_4\) 의 Niemeier 격자)이고, 그 근에 해당하는 \(W_{37}\) 의 노름 6 벡터 72쌍은 \(s=\pm9\), \(N(\nu)=\tfrac{37\cdot36}{16}-\tfrac{81}4=63<81=r^2\), \(\Gamma=\tfrac79\theta\).} 직접 인증(\texttt{certify.py}): \(O\cdot w\subset\Lambda\) (랭크 8) 의 행렬식이 예측식과 일치하고 \textbf{\(R_U\ge r^8D^4/\det(O\cdot w)=2.7326=(9/7)^4\)} --- \(O\) 의 극대성 가정 없이 성립. 이 불안정성은 같은 족의 모든 격자 \(\Lambda I\) (모든 ideal 류)로 전이된다(\(\det\) 비가 불변). 즉 \textbf{\(p=37\) 은 (b) 쪽: Weil 격자에 불균형한 짧은 벡터가 존재.} (11회차 휴리스틱의 예측 "위반 궤도 \(\approx3/[O^\times:O_F^\times]\) 개" 와 부합: 발견된 것은 \(703\times144=2|G|\) 개.)
- \textbf{{[}수치{]} \(p=61\) 은 지금까지의 모든 시험을 통과}(빌더 B): 분리 토러스 \(C_{15}\) 직선(랭크 8), \(C_5\)-고정(랭크 24), \(C_3\)-고정(랭크 40) 부분격자에서 \(t\)-격자 전체에 대한 전수 열거 --- 위반 없음, 각각 \(\min\Gamma=\theta\) 정확히. 랭크 120 전체에 대한 LLL 수준 탐색에서도 \(\theta\) 미만 없음(약한 증거). ⇒ \textbf{\(p=61\) (\(\dim W=120\), \(r=15\), 이득 \(g\approx63\)) 이 첫 미해결 사례.}
- 심사(교차검증) 진행 중 → 결과는 아래 마무리에 기록.

\textbf{다음 할 일 → 12회차}: \(p=61\) 에서 T2 의 필요조건(\(SL_2(5)\) icosian 성분의 \(e\ge3\))과 모든 군론적 구조의 \(\Delta\)-부분공간(부분군의 isotypic 성분)에 대해 \(R_U\) 를 체계적으로 계산.
